% answers_short.tex —— 2008 年·普通高等学校招生考试·天津卷·理科 · 简短答案速查
% 由 tools/make_short_answers.py 自 answers.tex 的 \gksolhead/\gkanswer 提取，
% 并与 2008/2008_数学真题_天津理_答案要点.md 逐题比对后生成，勿手改（改源 MD 后重跑）。

\begin{gknotice}
\textbf{使用说明}：本节只给最终答案便于快速核对，完整推导见解析版。答案为编写者独立推导，\textbf{非官方标准答案}。
\end{gknotice}

\gkpart{选择题}
\begin{center}\begin{tabular}{ccccccccccc}
\toprule
题号 & 1 & 2 & 3 & 4 & 5 & 6 & 7 & 8 & 9 & 10 \\
\midrule
答案 & A & D & B & C & B & A & D & C & A & B \\
\bottomrule
\end{tabular}\end{center}

\gkpart{填空题}
\begin{itemize}[leftmargin=2.6em]
  \item[\textbf{11.}] $40$。
  \item[\textbf{12.}] $24$。
  \item[\textbf{13.}] $x^2+(y-1)^2=10$。
  \item[\textbf{14.}] $3$。
  \item[\textbf{15.}] $\dfrac76$。
  \item[\textbf{16.}] $\{2\}$。
\end{itemize}

\gkpart{解答题}
\begin{itemize}[leftmargin=2.6em]
  \item[\textbf{17.}] （1）$\dfrac45$；（2）$-\dfrac{24+7\sqrt3}{50}$。
  \item[\textbf{18.}] （1）$p=\dfrac34$；（2）分布列见解析，$E\xi=2$。
  \item[\textbf{19.}] （1）证明见解析；（2）$\arccos\dfrac2{\sqrt{11}}$；（3）$\arccos\dfrac4{\sqrt{55}}$。
  \item[\textbf{20.}] （1）$a=-8$，$b=9$；（2）$a\leqslant0$ 时在 $(-\infty,0)$、$(0,+\infty)$ 上增，$a>0$ 时分区间讨论；（3）$b\leqslant\dfrac74$。
  \item[\textbf{21.}] （1）$\dfrac{x^2}4-\dfrac{y^2}5=1$；（2）$\left(-\infty,-\dfrac54\right)\cup\left(-\dfrac{\sqrt5}2,0\right)\cup\left(0,\dfrac{\sqrt5}2\right)\cup\left(\dfrac54,+\infty\right)$。
  \item[\textbf{22.}] （1）$a_2=3$，$b_2=9$；（2）$a_n=\dfrac{n(n+1)}2$，$b_n=(n+1)^2$；（3）证明见解析。
\end{itemize}
